\section*{Глава 1. Анализ предметной области}
\addcontentsline{toc}{section}{Глава 1. Анализ предметной области}
\stepcounter{section}

В данной главе выполняется анализ предметной области, связанной с хранением, систематизацией и публикацией учебных материалов в электронном виде. Рассматриваются существующие подходы и программные средства, формулируются цель, задачи и основные требования к разрабатываемому веб-приложению.

\subsection{Анализ существующих веб-приложений для хранения,\ систематизации и публикации учебных материалов}

В современных образовательных организациях учебные материалы представлены преимущественно в электронном виде. К таким материалам относятся лекции, лабораторные работы, методические указания, практические задания, презентации, шаблоны отчетов и иные вспомогательные документы. Все перечисленные данные должны быть доступны преподавателям и студентам в удобной и структурированной форме, а также поддерживаться в актуальном состоянии в течение учебного процесса.

На практике хранение учебных материалов часто организовано нецентрализованно. Документы размещаются в различных облачных сервисах, мессенджерах, на личных страницах преподавателей или в файловых каталогах. Подобный подход приводит к нескольким типичным проблемам. Во-первых, усложняется поиск нужного материала, так как пользователь вынужден обращаться к нескольким источникам. Во-вторых, отсутствует единая структура хранения, что затрудняет систематизацию материалов по дисциплинам, темам, видам занятий или авторам. В-третьих, при отсутствии ролевого разграничения доступа возрастает риск несанкционированного изменения или удаления данных. Наконец, становится затруднительным контроль актуальности материалов и учет действий пользователей.

Существующие программные решения частично закрывают обозначенные потребности. К наиболее распространенным подходам относятся облачные файловые хранилища, системы дистанционного обучения и внутренние корпоративные порталы кафедр и учебных подразделений.

\begin{center}
\begin{tabularx}{0.96\textwidth}{>{\raggedright\arraybackslash}p{3.8cm} >{\raggedright\arraybackslash}X >{\raggedright\arraybackslash}X}
\toprule
\textbf{Тип решения} & \textbf{Преимущества} & \textbf{Недостатки} \\
\midrule
Облачные файловые хранилища & Простая загрузка и скачивание файлов, совместный доступ, привычный интерфейс & Слабая предметно-ориентированная структура, ограниченные возможности ролевого разграничения и статистики \\
Системы дистанционного обучения & Поддержка учебного процесса, разграничение доступа, публикация материалов и заданий & Избыточность функциональности для узкой задачи электронного архива, более сложное сопровождение \\
Внутренние сайты и порталы кафедр & Возможность адаптации под локальные потребности, относительная простота публикации & Отсутствие единого подхода к моделированию данных, слабая поддержка CRUD-операций и аналитики \\
\bottomrule
\end{tabularx}
\end{center}

Облачные сервисы, такие как Google Drive и аналогичные решения, удобны для обмена файлами, однако ориентированы прежде всего на файловое хранение. В них отсутствует выраженная модель предметной области, которая необходима для учебного архива: категории материалов, авторство, ролевой доступ, фильтрация и агрегирование данных.

Системы дистанционного обучения, например Moodle и аналогичные платформы, предоставляют широкие возможности для организации учебного процесса. Они позволяют публиковать учебные материалы, управлять доступом и вести учет активности. Однако подобные решения во многих случаях являются слишком сложными для задач компактного архива учебных материалов. Для учебного проекта их архитектура и функциональная насыщенность могут быть избыточными.

Внутренние сайты кафедр и учебных подразделений обычно создаются под локальные задачи. Такие решения могут быть удобны в использовании, но часто развиваются бессистемно. В результате возникают трудности при расширении структуры данных, организации прав доступа и поддержке целостности информации.

Таким образом, проведенный анализ показывает целесообразность разработки отдельного веб-приложения, ориентированного именно на хранение, систематизацию и публикацию учебных материалов. Такое приложение должно быть проще полнофункциональных образовательных платформ, но при этом содержать обязательные механизмы аутентификации, авторизации, ролевого разграничения доступа, CRUD-операций, загрузки и скачивания файлов, а также базовой статистической обработки данных.

\subsection{Анализ программных средств разработки и выбор стека технологий}

Для реализации веб-приложения необходимо выбрать программные средства, позволяющие создать серверную логику, пользовательский интерфейс, механизм аутентификации, доступ к базе данных и обработку файлов. Поскольку проект выполняется в учебных целях, основными критериями выбора являются простота освоения, доступность документации, достаточность функциональности и удобство поэтапной разработки.

В качестве языка программирования выбран Python. Данный язык обладает простым и читаемым синтаксисом, широко используется при разработке веб-приложений и имеет большой набор библиотек для решения прикладных задач. В качестве альтернативы можно рассматривать Rust, который обеспечивает высокую производительность и безопасность памяти, однако его использование требует более глубокого понимания низкоуровневых механизмов и усложняет разработку учебного проекта. Поэтому для данной работы выбран Python, так как он позволяет быстрее реализовать основную логику приложения и сосредоточиться на предметной области.

В качестве веб-фреймворка выбран Flask. Он представляет собой легковесное средство разработки веб-приложений, позволяющее реализовать маршрутизацию, шаблонизацию страниц, обработку форм, работу с файлами и интеграцию со сторонними библиотеками. В качестве альтернативы можно использовать Django, который предоставляет больше встроенных возможностей, включая административную панель, ORM и готовую структуру проекта. Однако для небольшого учебного приложения Django может быть избыточным и требовать большего объема настройки. Flask выбран потому, что он проще по структуре и удобен для поэтапной реализации необходимых функций.

Для аутентификации пользователей и управления сессиями выбран модуль Flask-Login. Он предоставляет готовые механизмы входа в систему, выхода из нее, восстановления текущего пользователя и ограничения доступа к маршрутам для неавторизованных пользователей. В качестве альтернативы можно было реализовать механизм аутентификации самостоятельно с использованием сессий Flask. Однако самостоятельная реализация потребовала бы больше кода и увеличила бы вероятность ошибок при проверке авторизации. Flask-Login выбран потому, что он упрощает реализацию входа пользователей и делает работу с текущим пользователем более удобной и надежной.

Для работы с базой данных используется Flask-SQLAlchemy. Библиотека обеспечивает объектно-реляционное отображение, позволяя описывать сущности в виде Python-классов и выполнять операции с данными на уровне моделей. В качестве альтернативы можно использовать прямые SQL-запросы через модуль sqlite3. Такой подход дает больше контроля над запросами, однако усложняет поддержку кода и требует вручную описывать многие операции с данными. Flask-SQLAlchemy выбран потому, что он упрощает реализацию CRUD-интерфейса, уменьшает количество повторяющегося кода и делает структуру приложения более понятной.

В качестве системы управления базой данных выбрана SQLite. Ее преимуществами являются простота настройки, отсутствие необходимости устанавливать отдельный сервер базы данных и достаточность возможностей для учебного проекта. В качестве альтернативы можно использовать PostgreSQL, который лучше подходит для крупных и нагруженных приложений, поддерживает расширенные возможности и многопользовательскую работу. Однако для учебного проекта использование PostgreSQL усложнило бы настройку и развертывание. SQLite выбрана потому, что она не требует отдельного сервера и позволяет быстро приступить к разработке приложения.

Для построения пользовательского интерфейса используются HTML, CSS, Jinja2 и Bootstrap. HTML и CSS применяются для описания структуры и внешнего вида страниц, Jinja2 позволяет формировать динамические страницы на стороне сервера, а Bootstrap обеспечивает единый и аккуратный внешний вид интерфейса. В качестве альтернативы можно использовать отдельный JavaScript-фреймворк, например React. Такой подход удобен для сложных интерактивных интерфейсов, однако требует разработки отдельной клиентской части и дополнительной настройки взаимодействия с сервером. Для данного проекта выбран серверный рендеринг с Jinja2 и Bootstrap, так как он проще в реализации и полностью соответствует требованиям учебного веб-приложения.

\begin{center}
\footnotesize\begin{tabularx}{0.96\textwidth}{>{\raggedright\arraybackslash}p{4.1cm} >{\raggedright\arraybackslash}p{3.1cm} >{\raggedright\arraybackslash}X}
\toprule
\textbf{Компонент} & \textbf{Выбранное средство} & \textbf{Причина выбора} \\
\midrule
Язык программирования & Python & Простота синтаксиса, высокая скорость разработки, хорошая документация \\
Веб-фреймворк & Flask & Легковесность, гибкость, удобство для учебного проекта \\
Аутентификация & Flask-Login & Готовые механизмы входа, выхода и работы с пользовательской сессией \\
Доступ к БД & Flask-SQLAlchemy & Удобная работа с моделями, ORM-подход, упрощение CRUD-операций \\
СУБД & SQLite & Простая настройка, отсутствие отдельного серверного ПО, достаточность для проекта \\
Шаблоны и интерфейс & Jinja2, Bootstrap & Динамическая генерация страниц и единый визуальный стиль \\
Работа с файлами & Средства Flask и Werkzeug & Простая загрузка, сохранение и скачивание файлов \\
\bottomrule
\end{tabularx}
\end{center}

Таким образом, выбранный стек технологий включает Python, Flask, Flask-Login, Flask-SQLAlchemy, SQLite, Jinja2 и Bootstrap.

\subsection{Формулировка цели, задач и требований к системе}

Целью курсового проекта является разработка веб-приложения для хранения, систематизации и публикации учебных материалов с разграничением прав доступа пользователей в зависимости от их роли.

Для достижения поставленной цели необходимо решить следующие задачи:

\begin{itemize}[leftmargin=1.25cm]
    \item проанализировать предметную область хранения и предоставления учебных материалов;
    \item определить категории пользователей системы и их права доступа;
    \item выделить основные сущности предметной области и связи между ними;
    \item спроектировать структуру базы данных и интерфейс веб-приложения;
    \item реализовать аутентификацию и авторизацию пользователей;
    \item реализовать систему проверки прав доступа в соответствии с ролью пользователя;
    \item реализовать CRUD-интерфейс для пользователей, категорий и материалов;
    \item реализовать загрузку, скачивание и удаление файлов учебных материалов;
    \item реализовать агрегирование данных и формирование статистики;
    \item выполнить тестирование и подготовить пояснительную документацию.
\end{itemize}

В рамках системы выделяются следующие категории пользователей:

\begin{center}
\small\begin{tabularx}{0.96\textwidth}{>{\raggedright\arraybackslash}p{4.2cm} >{\raggedright\arraybackslash}X}
\toprule
\textbf{Роль} & \textbf{Основные функции} \\
\midrule
Администратор & Управление пользователями, категориями и всеми материалами, просмотр статистики, полный контроль над системой \\
Преподаватель & Создание и редактирование собственных материалов, загрузка и удаление файлов, просмотр категорий и материалов \\
Студент & Просмотр списка материалов, открытие карточек материалов, скачивание прикрепленных файлов \\
\bottomrule
\end{tabularx}
\end{center}

В приложении представлены следующие сущности:

\begin{center}
\begin{tabularx}{0.96\textwidth}{>{\raggedright\arraybackslash}p{3.4cm} >{\raggedright\arraybackslash}X}
\toprule
\textbf{Сущность} & \textbf{Назначение} \\
\midrule
Пользователь & Хранение сведений об учетной записи, логине, пароле, ФИО и роли \\
Роль & Определение уровня прав доступа пользователя \\
Категория & Классификация учебных материалов по тематике или типу \\
Материал & Основная сущность архива, содержащая название, описание, автора и категорию \\
Файл & Сведения о прикрепленном файле учебного материала и его хранении \\
\bottomrule
\end{tabularx}
\end{center}

К основным функциональным требованиям относятся:

\begin{itemize}[leftmargin=1.25cm]
    \item регистрация и аутентификация пользователей администраторами системы;
    \item авторизация и разграничение прав доступа по ролям;
    \item создание, просмотр, редактирование и удаление пользователей;
    \item создание, просмотр, редактирование и удаление категорий;
    \item создание, просмотр, редактирование и удаление материалов;
    \item загрузка, скачивание и удаление файлов, связанных с материалами;
    \item фильтрация и поиск материалов;
    \item формирование агрегированных данных и статистики.
\end{itemize}

К основным нефункциональным требованиям относятся удобство интерфейса, корректная работа в современных браузерах, поддержание целостности данных, ограничение доступа в соответствии с ролями и возможность дальнейшего расширения системы.